% Seminar 4 - Pagina 1. Fragment pentru
% % Seminar 4 - Pagina 1. Fragment pentru
% % Seminar 4 - Pagina 1. Fragment pentru
% % Seminar 4 - Pagina 1. Fragment pentru
% \input{seminars/pages/seminar4-pagina1-ro}; fără preambul sau pagină nouă.
% NOTE EDITORIALE:
% - Antetul este "ALGEBRA (Seminar 4)", iar data scrisă pare să fie
%   "31-OCT-95". Cronologia față de Seminarul 3 (24-oct.-1994) este
%   neclară; păstrăm anul 95, fără a-l corecta tacit.
% - La suma cardinalelor, manuscrisul pare să folosească A_i REUNIUNE
%   {i}, în loc de produsul cartezian A_i x {i}. Formula uzuală a
%   sumei folosește o reuniune disjunctă, de ex. |union_i(A_i x {i})|.
%   Păstrăm expresia manuscrisului.
% - Afirmația "operațiile sunt corect definite" este însoțită numai
%   de diagrama transferului unei aplicații B -> A prin bijecții.
% - Numerotările Ex. 4 și Ex. 5 apar de câte două ori; le păstrăm.
% - Paradoxul lui Cantor este formulat în termenii "mulțimii tuturor
%   cardinalelor". Distincția modernă de clasă proprie nu este inserată
%   în conținutul manuscrisului.
% - Demonstrația pentru majorarea unei mulțimi de cardinali folosește
%   formula sumei menționate mai sus, care cere o reuniune disjunctă.
% - Pasajele barate și urmele foarte slabe nu sunt reconstruite.

\phantomsection
\subsection*{Seminar 4}
\addcontentsline{toc}{subsection}{Seminar 4}
\begin{flushright}
\textit{31-oct.-1995 
% (conform manuscrisului)
}
\end{flushright}
\begingroup
\setlength{\abovedisplayskip}{5pt}
\setlength{\belowdisplayskip}{5pt}
\setlength{\abovedisplayshortskip}{3pt}
\setlength{\belowdisplayshortskip}{3pt}

\noindent\underline{Operații cu numere cardinale.}
Fie $\alpha=|A|$, $\beta=|B|$. Definim
\[
  \alpha\beta\ \overset{\mathrm{def}}{=}\ |A\times B|,
  \qquad
  \alpha^{\beta}\ \overset{\mathrm{def}}{=}\
     |\operatorname{Hom}(B,A)|.
\]
Pentru $\alpha=2$,
\[
  2^{\beta}=|\mathcal P(B)|.
\]
Fie $(\alpha_i)_{i\in I}$ o familie de numere cardinale,
$\alpha_i=|A_i|$. În manuscris suma este definită prin
\[
  \sum_{i\in I}\alpha_i\ \overset{\mathrm{def}}{=}\
  \left|\bigcup_{i\in I}(A_i\cup\{i\})\right|.
\]

\noindent\underline{Operațiile sunt corect definite.}
Presupunem că
\[
\begin{aligned}
 \alpha=|A|=|A'|&\iff\exists\,f:A\longrightarrow A'\text{ bijectivă},\\
 \beta=|B|=|B'|&\iff\exists\,g:B\longrightarrow B'\text{ bijectivă}.
\end{aligned}
\]
Pentru orice $h:B\longrightarrow A$, se obține
\[
  \varphi(h)=f\circ h\circ g^{-1},\qquad\varphi\text{ bijectivă}.
\]
\[
\begin{tikzcd}[column sep=large,row sep=large]
 B\arrow[r,"h"] & A\arrow[d,"f"]\\
 B'\arrow[u,"g^{-1}"]\arrow[r,"f\circ h\circ g^{-1}"'] & A'
\end{tikzcd}
\]

\medskip
\noindent\underline{Exercițiul 4.}
Fie $K$ o mulțime de numere cardinale. Să se arate că există un
număr cardinal mai mare decât orice cardinal din $K$.

\smallskip
\noindent\underline{Exercițiul 5 (paradoxul lui Cantor).}
Mulțimea tuturor numerelor cardinale duce la contradicție.

\smallskip
\noindent\underline{Exercițiul 4 (numerotare repetată în manuscris).}
Pentru $\alpha=|A|$ și $\beta=|B|$, definițiile notate sunt
\[
 \alpha\leq\beta\ \overset{\mathrm{def}}{\iff}\
 \exists\,f:A\longrightarrow B\text{ injectivă},
 \qquad
 \alpha<\beta\ \overset{\mathrm{def}}{\iff}\
 \alpha\leq\beta\ \text{și}\ \alpha\neq\beta.
\]
\noindent\underline{Observație.}
Nu există aplicație surjectivă $A\longrightarrow\mathcal P(A)$,
iar $h:A\longrightarrow\mathcal P(A)$, $h(a)=\{a\}$, este injectivă.
Prin urmare,
\[
 \alpha=|A|<|\mathcal P(A)|\iff\alpha<2^{\alpha}.
\]

\noindent\underline{Demonstrația primului exercițiu 4.}
Scriem $K=\{\alpha_i\mid i\in I\}$ și notăm
$\beta\overset{\mathrm{not}}{=}\sum_{i\in I}\alpha_i$.
În manuscris se folosesc aplicațiile
\[
 \varphi_i:A_i\longrightarrow\bigcup_{i\in I}(A_i\cup\{i\}),
 \qquad\varphi_i(x)=x,
\]
și rezultă $\alpha_i\leq\beta$ pentru orice $i\in I$.
Cum $\beta<2^{\beta}$, notăm $\gamma=2^{\beta}$. Atunci
$\alpha_i<\gamma$ pentru orice $i\in I$.

\noindent\underline{Exercițiul 5 (reluat).}
Dacă $\mathcal C$ este mulțimea tuturor numerelor cardinale,
argumentul precedent dă $\gamma\in\mathcal C$ cu
$\gamma>\alpha$ pentru orice $\alpha\in\mathcal C$,
ceea ce este o contradicție.
\endgroup
; fără preambul sau pagină nouă.
% NOTE EDITORIALE:
% - Antetul este "ALGEBRA (Seminar 4)", iar data scrisă pare să fie
%   "31-OCT-95". Cronologia față de Seminarul 3 (24-oct.-1994) este
%   neclară; păstrăm anul 95, fără a-l corecta tacit.
% - La suma cardinalelor, manuscrisul pare să folosească A_i REUNIUNE
%   {i}, în loc de produsul cartezian A_i x {i}. Formula uzuală a
%   sumei folosește o reuniune disjunctă, de ex. |union_i(A_i x {i})|.
%   Păstrăm expresia manuscrisului.
% - Afirmația "operațiile sunt corect definite" este însoțită numai
%   de diagrama transferului unei aplicații B -> A prin bijecții.
% - Numerotările Ex. 4 și Ex. 5 apar de câte două ori; le păstrăm.
% - Paradoxul lui Cantor este formulat în termenii "mulțimii tuturor
%   cardinalelor". Distincția modernă de clasă proprie nu este inserată
%   în conținutul manuscrisului.
% - Demonstrația pentru majorarea unei mulțimi de cardinali folosește
%   formula sumei menționate mai sus, care cere o reuniune disjunctă.
% - Pasajele barate și urmele foarte slabe nu sunt reconstruite.

\phantomsection
\subsection*{Seminar 4}
\addcontentsline{toc}{subsection}{Seminar 4}
\begin{flushright}
\textit{31-oct.-1995 
% (conform manuscrisului)
}
\end{flushright}
\begingroup
\setlength{\abovedisplayskip}{5pt}
\setlength{\belowdisplayskip}{5pt}
\setlength{\abovedisplayshortskip}{3pt}
\setlength{\belowdisplayshortskip}{3pt}

\noindent\underline{Operații cu numere cardinale.}
Fie $\alpha=|A|$, $\beta=|B|$. Definim
\[
  \alpha\beta\ \overset{\mathrm{def}}{=}\ |A\times B|,
  \qquad
  \alpha^{\beta}\ \overset{\mathrm{def}}{=}\
     |\operatorname{Hom}(B,A)|.
\]
Pentru $\alpha=2$,
\[
  2^{\beta}=|\mathcal P(B)|.
\]
Fie $(\alpha_i)_{i\in I}$ o familie de numere cardinale,
$\alpha_i=|A_i|$. În manuscris suma este definită prin
\[
  \sum_{i\in I}\alpha_i\ \overset{\mathrm{def}}{=}\
  \left|\bigcup_{i\in I}(A_i\cup\{i\})\right|.
\]

\noindent\underline{Operațiile sunt corect definite.}
Presupunem că
\[
\begin{aligned}
 \alpha=|A|=|A'|&\iff\exists\,f:A\longrightarrow A'\text{ bijectivă},\\
 \beta=|B|=|B'|&\iff\exists\,g:B\longrightarrow B'\text{ bijectivă}.
\end{aligned}
\]
Pentru orice $h:B\longrightarrow A$, se obține
\[
  \varphi(h)=f\circ h\circ g^{-1},\qquad\varphi\text{ bijectivă}.
\]
\[
\begin{tikzcd}[column sep=large,row sep=large]
 B\arrow[r,"h"] & A\arrow[d,"f"]\\
 B'\arrow[u,"g^{-1}"]\arrow[r,"f\circ h\circ g^{-1}"'] & A'
\end{tikzcd}
\]

\medskip
\noindent\underline{Exercițiul 4.}
Fie $K$ o mulțime de numere cardinale. Să se arate că există un
număr cardinal mai mare decât orice cardinal din $K$.

\smallskip
\noindent\underline{Exercițiul 5 (paradoxul lui Cantor).}
Mulțimea tuturor numerelor cardinale duce la contradicție.

\smallskip
\noindent\underline{Exercițiul 4 (numerotare repetată în manuscris).}
Pentru $\alpha=|A|$ și $\beta=|B|$, definițiile notate sunt
\[
 \alpha\leq\beta\ \overset{\mathrm{def}}{\iff}\
 \exists\,f:A\longrightarrow B\text{ injectivă},
 \qquad
 \alpha<\beta\ \overset{\mathrm{def}}{\iff}\
 \alpha\leq\beta\ \text{și}\ \alpha\neq\beta.
\]
\noindent\underline{Observație.}
Nu există aplicație surjectivă $A\longrightarrow\mathcal P(A)$,
iar $h:A\longrightarrow\mathcal P(A)$, $h(a)=\{a\}$, este injectivă.
Prin urmare,
\[
 \alpha=|A|<|\mathcal P(A)|\iff\alpha<2^{\alpha}.
\]

\noindent\underline{Demonstrația primului exercițiu 4.}
Scriem $K=\{\alpha_i\mid i\in I\}$ și notăm
$\beta\overset{\mathrm{not}}{=}\sum_{i\in I}\alpha_i$.
În manuscris se folosesc aplicațiile
\[
 \varphi_i:A_i\longrightarrow\bigcup_{i\in I}(A_i\cup\{i\}),
 \qquad\varphi_i(x)=x,
\]
și rezultă $\alpha_i\leq\beta$ pentru orice $i\in I$.
Cum $\beta<2^{\beta}$, notăm $\gamma=2^{\beta}$. Atunci
$\alpha_i<\gamma$ pentru orice $i\in I$.

\noindent\underline{Exercițiul 5 (reluat).}
Dacă $\mathcal C$ este mulțimea tuturor numerelor cardinale,
argumentul precedent dă $\gamma\in\mathcal C$ cu
$\gamma>\alpha$ pentru orice $\alpha\in\mathcal C$,
ceea ce este o contradicție.
\endgroup
; fără preambul sau pagină nouă.
% NOTE EDITORIALE:
% - Antetul este "ALGEBRA (Seminar 4)", iar data scrisă pare să fie
%   "31-OCT-95". Cronologia față de Seminarul 3 (24-oct.-1994) este
%   neclară; păstrăm anul 95, fără a-l corecta tacit.
% - La suma cardinalelor, manuscrisul pare să folosească A_i REUNIUNE
%   {i}, în loc de produsul cartezian A_i x {i}. Formula uzuală a
%   sumei folosește o reuniune disjunctă, de ex. |union_i(A_i x {i})|.
%   Păstrăm expresia manuscrisului.
% - Afirmația "operațiile sunt corect definite" este însoțită numai
%   de diagrama transferului unei aplicații B -> A prin bijecții.
% - Numerotările Ex. 4 și Ex. 5 apar de câte două ori; le păstrăm.
% - Paradoxul lui Cantor este formulat în termenii "mulțimii tuturor
%   cardinalelor". Distincția modernă de clasă proprie nu este inserată
%   în conținutul manuscrisului.
% - Demonstrația pentru majorarea unei mulțimi de cardinali folosește
%   formula sumei menționate mai sus, care cere o reuniune disjunctă.
% - Pasajele barate și urmele foarte slabe nu sunt reconstruite.

\phantomsection
\subsection*{Seminar 4}
\addcontentsline{toc}{subsection}{Seminar 4}
\begin{flushright}
\textit{31-oct.-1995 
% (conform manuscrisului)
}
\end{flushright}
\begingroup
\setlength{\abovedisplayskip}{5pt}
\setlength{\belowdisplayskip}{5pt}
\setlength{\abovedisplayshortskip}{3pt}
\setlength{\belowdisplayshortskip}{3pt}

\noindent\underline{Operații cu numere cardinale.}
Fie $\alpha=|A|$, $\beta=|B|$. Definim
\[
  \alpha\beta\ \overset{\mathrm{def}}{=}\ |A\times B|,
  \qquad
  \alpha^{\beta}\ \overset{\mathrm{def}}{=}\
     |\operatorname{Hom}(B,A)|.
\]
Pentru $\alpha=2$,
\[
  2^{\beta}=|\mathcal P(B)|.
\]
Fie $(\alpha_i)_{i\in I}$ o familie de numere cardinale,
$\alpha_i=|A_i|$. În manuscris suma este definită prin
\[
  \sum_{i\in I}\alpha_i\ \overset{\mathrm{def}}{=}\
  \left|\bigcup_{i\in I}(A_i\cup\{i\})\right|.
\]

\noindent\underline{Operațiile sunt corect definite.}
Presupunem că
\[
\begin{aligned}
 \alpha=|A|=|A'|&\iff\exists\,f:A\longrightarrow A'\text{ bijectivă},\\
 \beta=|B|=|B'|&\iff\exists\,g:B\longrightarrow B'\text{ bijectivă}.
\end{aligned}
\]
Pentru orice $h:B\longrightarrow A$, se obține
\[
  \varphi(h)=f\circ h\circ g^{-1},\qquad\varphi\text{ bijectivă}.
\]
\[
\begin{tikzcd}[column sep=large,row sep=large]
 B\arrow[r,"h"] & A\arrow[d,"f"]\\
 B'\arrow[u,"g^{-1}"]\arrow[r,"f\circ h\circ g^{-1}"'] & A'
\end{tikzcd}
\]

\medskip
\noindent\underline{Exercițiul 4.}
Fie $K$ o mulțime de numere cardinale. Să se arate că există un
număr cardinal mai mare decât orice cardinal din $K$.

\smallskip
\noindent\underline{Exercițiul 5 (paradoxul lui Cantor).}
Mulțimea tuturor numerelor cardinale duce la contradicție.

\smallskip
\noindent\underline{Exercițiul 4 (numerotare repetată în manuscris).}
Pentru $\alpha=|A|$ și $\beta=|B|$, definițiile notate sunt
\[
 \alpha\leq\beta\ \overset{\mathrm{def}}{\iff}\
 \exists\,f:A\longrightarrow B\text{ injectivă},
 \qquad
 \alpha<\beta\ \overset{\mathrm{def}}{\iff}\
 \alpha\leq\beta\ \text{și}\ \alpha\neq\beta.
\]
\noindent\underline{Observație.}
Nu există aplicație surjectivă $A\longrightarrow\mathcal P(A)$,
iar $h:A\longrightarrow\mathcal P(A)$, $h(a)=\{a\}$, este injectivă.
Prin urmare,
\[
 \alpha=|A|<|\mathcal P(A)|\iff\alpha<2^{\alpha}.
\]

\noindent\underline{Demonstrația primului exercițiu 4.}
Scriem $K=\{\alpha_i\mid i\in I\}$ și notăm
$\beta\overset{\mathrm{not}}{=}\sum_{i\in I}\alpha_i$.
În manuscris se folosesc aplicațiile
\[
 \varphi_i:A_i\longrightarrow\bigcup_{i\in I}(A_i\cup\{i\}),
 \qquad\varphi_i(x)=x,
\]
și rezultă $\alpha_i\leq\beta$ pentru orice $i\in I$.
Cum $\beta<2^{\beta}$, notăm $\gamma=2^{\beta}$. Atunci
$\alpha_i<\gamma$ pentru orice $i\in I$.

\noindent\underline{Exercițiul 5 (reluat).}
Dacă $\mathcal C$ este mulțimea tuturor numerelor cardinale,
argumentul precedent dă $\gamma\in\mathcal C$ cu
$\gamma>\alpha$ pentru orice $\alpha\in\mathcal C$,
ceea ce este o contradicție.
\endgroup
; fără preambul sau pagină nouă.
% NOTE EDITORIALE:
% - Antetul este "ALGEBRA (Seminar 4)", iar data scrisă pare să fie
%   "31-OCT-95". Cronologia față de Seminarul 3 (24-oct.-1994) este
%   neclară; păstrăm anul 95, fără a-l corecta tacit.
% - La suma cardinalelor, manuscrisul pare să folosească A_i REUNIUNE
%   {i}, în loc de produsul cartezian A_i x {i}. Formula uzuală a
%   sumei folosește o reuniune disjunctă, de ex. |union_i(A_i x {i})|.
%   Păstrăm expresia manuscrisului.
% - Afirmația "operațiile sunt corect definite" este însoțită numai
%   de diagrama transferului unei aplicații B -> A prin bijecții.
% - Numerotările Ex. 4 și Ex. 5 apar de câte două ori; le păstrăm.
% - Paradoxul lui Cantor este formulat în termenii "mulțimii tuturor
%   cardinalelor". Distincția modernă de clasă proprie nu este inserată
%   în conținutul manuscrisului.
% - Demonstrația pentru majorarea unei mulțimi de cardinali folosește
%   formula sumei menționate mai sus, care cere o reuniune disjunctă.
% - Pasajele barate și urmele foarte slabe nu sunt reconstruite.

\phantomsection
\subsection*{Seminar 4}
\addcontentsline{toc}{subsection}{Seminar 4}
\begin{flushright}
\textit{31-oct.-1995 
% (conform manuscrisului)
}
\end{flushright}
\begingroup
\setlength{\abovedisplayskip}{5pt}
\setlength{\belowdisplayskip}{5pt}
\setlength{\abovedisplayshortskip}{3pt}
\setlength{\belowdisplayshortskip}{3pt}

\noindent\underline{Operații cu numere cardinale.}
Fie $\alpha=|A|$, $\beta=|B|$. Definim
\[
  \alpha\beta\ \overset{\mathrm{def}}{=}\ |A\times B|,
  \qquad
  \alpha^{\beta}\ \overset{\mathrm{def}}{=}\
     |\operatorname{Hom}(B,A)|.
\]
Pentru $\alpha=2$,
\[
  2^{\beta}=|\mathcal P(B)|.
\]
Fie $(\alpha_i)_{i\in I}$ o familie de numere cardinale,
$\alpha_i=|A_i|$. În manuscris suma este definită prin
\[
  \sum_{i\in I}\alpha_i\ \overset{\mathrm{def}}{=}\
  \left|\bigcup_{i\in I}(A_i\cup\{i\})\right|.
\]

\noindent\underline{Operațiile sunt corect definite.}
Presupunem că
\[
\begin{aligned}
 \alpha=|A|=|A'|&\iff\exists\,f:A\longrightarrow A'\text{ bijectivă},\\
 \beta=|B|=|B'|&\iff\exists\,g:B\longrightarrow B'\text{ bijectivă}.
\end{aligned}
\]
Pentru orice $h:B\longrightarrow A$, se obține
\[
  \varphi(h)=f\circ h\circ g^{-1},\qquad\varphi\text{ bijectivă}.
\]
\[
\begin{tikzcd}[column sep=large,row sep=large]
 B\arrow[r,"h"] & A\arrow[d,"f"]\\
 B'\arrow[u,"g^{-1}"]\arrow[r,"f\circ h\circ g^{-1}"'] & A'
\end{tikzcd}
\]

\medskip
\noindent\underline{Exercițiul 4.}
Fie $K$ o mulțime de numere cardinale. Să se arate că există un
număr cardinal mai mare decât orice cardinal din $K$.

\smallskip
\noindent\underline{Exercițiul 5 (paradoxul lui Cantor).}
Mulțimea tuturor numerelor cardinale duce la contradicție.

\smallskip
\noindent\underline{Exercițiul 4 (numerotare repetată în manuscris).}
Pentru $\alpha=|A|$ și $\beta=|B|$, definițiile notate sunt
\[
 \alpha\leq\beta\ \overset{\mathrm{def}}{\iff}\
 \exists\,f:A\longrightarrow B\text{ injectivă},
 \qquad
 \alpha<\beta\ \overset{\mathrm{def}}{\iff}\
 \alpha\leq\beta\ \text{și}\ \alpha\neq\beta.
\]
\noindent\underline{Observație.}
Nu există aplicație surjectivă $A\longrightarrow\mathcal P(A)$,
iar $h:A\longrightarrow\mathcal P(A)$, $h(a)=\{a\}$, este injectivă.
Prin urmare,
\[
 \alpha=|A|<|\mathcal P(A)|\iff\alpha<2^{\alpha}.
\]

\noindent\underline{Demonstrația primului exercițiu 4.}
Scriem $K=\{\alpha_i\mid i\in I\}$ și notăm
$\beta\overset{\mathrm{not}}{=}\sum_{i\in I}\alpha_i$.
În manuscris se folosesc aplicațiile
\[
 \varphi_i:A_i\longrightarrow\bigcup_{i\in I}(A_i\cup\{i\}),
 \qquad\varphi_i(x)=x,
\]
și rezultă $\alpha_i\leq\beta$ pentru orice $i\in I$.
Cum $\beta<2^{\beta}$, notăm $\gamma=2^{\beta}$. Atunci
$\alpha_i<\gamma$ pentru orice $i\in I$.

\noindent\underline{Exercițiul 5 (reluat).}
Dacă $\mathcal C$ este mulțimea tuturor numerelor cardinale,
argumentul precedent dă $\gamma\in\mathcal C$ cu
$\gamma>\alpha$ pentru orice $\alpha\in\mathcal C$,
ceea ce este o contradicție.
\endgroup
